\paragraph{零温 witness 复杂度、Ramanujan 谱缝与平衡均值间隙}
在输出位势的端点大偏差、地基熵、局部 Legendre 几何与碰撞位势的 RH/Ramanujan 临界窗已经建立之后，还可进一步把零温极端层的可恢复性、可采样性与谱窗口失效压缩为下列若干条结果。统一记
\[
\mathcal G_n:=
\{\gamma:\ |\gamma|=n,\ \text{其输出 }1\text{ 的总数为 }k_{\max}(n)\}
\]
为极端饱和层。

\begin{theorem}[极端饱和层的完备外置阈值恰为 \texorpdfstring{$\log_2 c$}{log2 c} 比特/步]\label{thm:xi-real-input-ground-layer-oracle-threshold}
\leanverified{paper\_xi\_real\_input\_ground\_layer\_oracle\_threshold}
沿用推论 \ref{cor:real-input-40-ground-entropy} 的记号，并令 \(\Gamma_n\) 在 \(\mathcal G_n\) 上均匀分布。对任意只允许 \(B_n\) 比特外置寄存器的编码--解码机制，记恢复输出为 \(\widehat\Gamma_n\)。则有
\[
\boxed{\;
\mathbb P(\widehat\Gamma_n=\Gamma_n)
\le
\frac{2^{B_n}}{A_{n,k_{\max}(n)}}.
\;}
\]
从而：
\begin{enumerate}
  \item 若存在 \(\varepsilon>0\) 使
  \[
  B_n\le (\log_2 c-\varepsilon)n
  \]
  对充分大的 \(n\) 成立，则
  \[
  \mathbb P(\widehat\Gamma_n=\Gamma_n)
  \le
  \exp\!\bigl(-\varepsilon(\log 2)\,n+o(n)\bigr).
  \]
  \item 若希望存在固定 \(\eta\in(0,1]\) 使
  \[
  \mathbb P(\widehat\Gamma_n=\Gamma_n)\ge \eta
  \]
  对充分大的 \(n\) 成立，则必有
  \[
  B_n\ge n\log_2 c+\log_2\eta+o(n).
  \]
\end{enumerate}
换言之，极端饱和层的最小完备外置率恰为 \(\log_2 c\) 比特/步。
\end{theorem}
\begin{proof}
由定义，
\[
|\mathcal G_n|=A_{n,k_{\max}(n)}.
\]
设外置寄存器的取值集合为 \(\mathcal R_n\)，则
\(
|\mathcal R_n|\le 2^{B_n}
\)。
对任意确定性编码--解码机制，每个寄存器值 \(r\in\mathcal R_n\) 至多只能对应一个被正确恢复的候选路径；因此在 \(\mathcal G_n\) 中至多只有 \(2^{B_n}\) 条路径可能被正确译回。由于 \(\Gamma_n\) 在 \(\mathcal G_n\) 上均匀分布，得到
\[
\mathbb P(\widehat\Gamma_n=\Gamma_n)\le \frac{2^{B_n}}{|\mathcal G_n|}
=\frac{2^{B_n}}{A_{n,k_{\max}(n)}}.
\]
随机机制只是确定性机制的凸组合，故同一上界仍然成立。

再由推论 \ref{cor:real-input-40-ground-entropy}，
\[
A_{n,k_{\max}(n)}=\exp\bigl((\log c)n+o(n)\bigr).
\]
将其代回上式，便得到
\[
\mathbb P(\widehat\Gamma_n=\Gamma_n)
\le
\exp\bigl(B_n\log 2-(\log c)n+o(n)\bigr).
\]
若 \(B_n\le (\log_2 c-\varepsilon)n\)，则
\[
B_n\log 2-(\log c)n\le -\varepsilon(\log 2)\,n,
\]
从而得到第一条。
反之，若成功概率对充分大的 \(n\) 满足
\(
\mathbb P(\widehat\Gamma_n=\Gamma_n)\ge \eta
\)，则必有
\[
\eta
\le
\frac{2^{B_n}}{A_{n,k_{\max}(n)}},
\]
即
\[
B_n\ge \log_2 A_{n,k_{\max}(n)}+\log_2\eta
=n\log_2 c+\log_2\eta+o(n).
\]
证毕。
\end{proof}

\begin{theorem}[极端饱和层在无偏路径总体中的稀有指数等于 \texorpdfstring{$\log(\varphi^2/c)$}{log(phi^2/c)}]\label{thm:xi-real-input-ground-layer-unbiased-rarity}
\leanverified{paper\_xi\_real\_input\_ground\_layer\_unbiased\_rarity}
沿用命题 \ref{prop:real-input-40-pressure-freezing} 与推论 \ref{cor:real-input-40-ground-entropy} 的记号，记
\[
Z_n(1):=\mathbf{1}^{\mathsf T}M(1)^n\mathbf{1},
\qquad
p_n:=\frac{A_{n,k_{\max}(n)}}{Z_n(1)}.
\]
则有精确指数律
\[
\boxed{\;
\lim_{n\to\infty}\frac{1}{n}\log p_n
=
\log c-\log\lambda(1)
=
-\delta_\ast,
\;}
\]
其中
\[
\lambda(1)=\frac{3+\sqrt5}{2}=\varphi^2,
\qquad
\delta_\ast:=\log\frac{\varphi^2}{c}\approx 0.32249108560171.
\]
特别地，
\[
p_n=\exp(-\delta_\ast n+o(n)).
\]
若在全部长度 \(n\) 的可行路径上做独立无偏抽样，直到首次命中 \(\mathcal G_n\) 的等待时间记为 \(T_n\)，则
\[
\mathbb E[T_n]=\frac{1}{p_n}=\exp(\delta_\ast n+o(n)).
\]
\end{theorem}
\begin{proof}
由命题 \ref{prop:real-input-40-pressure-freezing}，
\[
\lim_{n\to\infty}\frac{1}{n}\log Z_n(1)=P(0)=\log\lambda(1)=\log(\varphi^2).
\]
另一方面，推论 \ref{cor:real-input-40-ground-entropy} 给出
\[
\lim_{n\to\infty}\frac{1}{n}\log A_{n,k_{\max}(n)}=\log c.
\]
两式相减即得
\[
\lim_{n\to\infty}\frac{1}{n}\log p_n
=
\log c-\log(\varphi^2)
=
-\delta_\ast.
\]
这与推论 \ref{cor:real-input-40-rate-endpoints} 的端点代价
\(
I(1/2)=\log(\varphi^2/c)
\)
完全一致。
若独立重复无偏抽样，则 \(T_n\) 服从参数为 \(p_n\) 的几何分布，因而
\(
\mathbb E[T_n]=1/p_n
\)，从而得到所示指数律。证毕。
\end{proof}

\begin{theorem}[RH-sharp 临界参数严格位于零温四次锚点之前]\label{thm:xi-real-input-rh-window-zero-temp-perron-gap}
\leanverified{paper\_xi\_real\_input\_rh\_window\_zero\_temp\_perron\_gap}
沿用命题 \ref{prop:real-input-40-collision-rhk-window} 与定理 \ref{thm:xi-zero-temperature-fourstate-artin-mazur-zeta-primitive-asymptotic} 的记号。设 \(u_R>1\) 为唯一满足
\[
u_R^5+4u_R^4+3u_R^3-96u_R^2-42u_R-14=0
\]
的正实根，并记
\[
\mathcal Q(u):=u^5+4u^4+3u^3-96u^2-42u-14,
\]
\[
\mathcal R(y):=y^4-6y^3+9y^2-y-1,
\qquad
y_\infty:=\rho(B_\infty)=c^2.
\]
则有严格有理夹逼
\[
\boxed{\;
\frac{359}{100}<u_R<\frac{3593}{1000}<y_\infty<\frac{18}{5}.\;}
\]
特别地，
\[
\boxed{\;
u_R<y_\infty=\rho(B_\infty)=c^2.
\;}
\]
因此在严格非空区间
\[
u\in\bigl(u_R,\rho(B_\infty)\bigr)
\]
上，核版 RH/Ramanujan 条件已经失效，但零温 quartic 地基尺度尚未达到其自身 Perron 阈值。
\end{theorem}
\begin{proof}
由命题 \ref{prop:real-input-40-collision-rhk-window}，\(\mathcal Q\) 在 \((0,\infty)\) 上恰有一个正根 \(u_R\)。直接代入得到
\[
\mathcal Q\!\left(\frac{359}{100}\right)
=
-\frac{25100699801}{10^{10}}<0,
\qquad
\mathcal Q\!\left(\frac{3593}{1000}\right)
=
\frac{362872814205193}{10^{15}}>0.
\]
由介值定理与唯一正根性，立得
\[
\frac{359}{100}<u_R<\frac{3593}{1000}.
\]

另一方面，定理 \ref{thm:xi-zero-temperature-fourstate-artin-mazur-zeta-primitive-asymptotic} 给出 \(y_\infty\) 是四次方程
\[
\mathcal R(y)=y^4-6y^3+9y^2-y-1=0
\]
的最大正根。
\[
\mathcal R\!\left(\frac{3593}{1000}\right)
=
-\frac{53334838799}{10^{12}}<0,
\qquad
\mathcal R\!\left(\frac{18}{5}\right)
=
\frac{41}{625}>0.
\]
因此同样由介值定理可得
\[
\frac{3593}{1000}<y_\infty<\frac{18}{5}.
\]
\[
u_R<y_\infty=\rho(B_\infty)=c^2.
\]
\[
u_R<\frac{3593}{1000}<y_\infty=\rho(B_\infty)=c^2.
\]
于是所述严格非空区间
\[
u\in\bigl(u_R,\rho(B_\infty)\bigr)
\]
立刻成立。证毕。
\end{proof}

\begin{theorem}[RH-sharp 阈与零温冻结阈的代数分裂]\label{thm:xi-real-input-rh-threshold-zero-temp-algebraic-splitting}
\leanverified{paper\_xi\_real\_input\_rh\_threshold\_zero\_temp\_algebraic\_splitting}
沿用定理 \ref{thm:xi-real-input-rh-window-zero-temp-perron-gap} 的记号。则
\[
\gcd_{\ZZ[X]}\bigl(\mathcal Q(X),\mathcal R(X)\bigr)=1.
\]
等价地，
\[
\operatorname{Res}\bigl(\mathcal Q,\mathcal R\bigr)\neq 0.
\]
因此 \(u_R\) 与 \(y_\infty=\rho(B_\infty)\) 不可能由同一个代数根同时满足这两条定义方程；特别地，
\[
u_R\neq y_\infty.
\]
结合定理 \ref{thm:xi-real-input-rh-window-zero-temp-perron-gap}，可知 RH-sharp 谱阈与零温冻结阈不仅在实轴上严格分离，而且在代数上亦不可合并。
\end{theorem}
\begin{proof}
对整系数多项式 \(\mathcal Q\) 与 \(\mathcal R\) 施行 Euclid 算法可得
\[
\gcd_{\ZZ[X]}\bigl(\mathcal Q(X),\mathcal R(X)\bigr)=1,
\]
等价地，
\[
\operatorname{Res}\bigl(\mathcal Q,\mathcal R\bigr)=-720892\neq 0.
\]
故 \(\mathcal Q\) 与 \(\mathcal R\) 没有公共代数根。

另一方面，定理 \ref{thm:xi-real-input-rh-window-zero-temp-perron-gap} 已给出
\[
\mathcal Q(u_R)=0,
\qquad
\mathcal R(y_\infty)=0,
\qquad
u_R<y_\infty.
\]
若 \(u_R=y_\infty\)，则该代数数将同时是 \(\mathcal Q\) 与 \(\mathcal R\) 的公共根，矛盾。故
\[
u_R\neq y_\infty.
\]
于是两条阈值边界在定义方程层面已经代数分裂，而定理 \ref{thm:xi-real-input-rh-window-zero-temp-perron-gap} 又进一步给出其实轴顺序为
\(
u_R<y_\infty
\)。
证毕。
\end{proof}

\endinput
\paragraph{输出位势速率函数的端点对数尖点与不对称刚性}
在定理 \ref{thm:xi-real-input-rh-window-zero-temp-perron-gap} 已把 RH-sharp 五次阈值与零温四次锚点严格分离、定理 \ref{thm:xi-output-potential-zero-temp-square-root-critical-law} 已锁定右端零温尾部的 \(e^{-\theta/2}\) 量级之后，归一化速率函数本身在两端还呈现出更细且彼此不对称的对数边界几何。

\begin{theorem}[输出位势速率函数在右端点具有显式对数尖点]\label{thm:xi-output-potential-rate-right-endpoint-log-cusp}\leanverified{paper\_xi\_output\_potential\_rate\_right\_endpoint\_log\_cusp}
沿用命题 \ref{prop:real-input-40-pressure-freezing}、推论 \ref{cor:real-input-40-rate-endpoints} 与结论 \ref{con:xi-terminal-real-input-output-microcanonical-logcusp-second-order} 的记号。记
\[
I(\alpha):=\sup_{\theta\in\RR}\bigl(\theta\alpha-(P(\theta)-P(0))\bigr),
\qquad
A:=\frac{d}{2c}>0,
\]
并令
\[
\varepsilon:=\frac12-\alpha.
\]
则当 \(\alpha\uparrow \frac12\)（即 \(\varepsilon\downarrow 0\)）时，有
\[
\boxed{\;
I(\alpha)
=
I\!\left(\frac12\right)
+2\varepsilon\log\varepsilon
-2(1+\log A)\varepsilon
+O(\varepsilon^2).\;}
\]
等价地，
\[
\boxed{\;
I\!\left(\frac12\right)-I(\alpha)
=
2\varepsilon|\log\varepsilon|
+2(1+\log A)\varepsilon
+O(\varepsilon^2).\;}
\]
并且导数满足精确发散律
\[
\boxed{\;
I'(\alpha)
=
-2\log\!\Bigl(\frac12-\alpha\Bigr)+2\log A+O\!\Bigl(\frac12-\alpha\Bigr),
\qquad
\alpha\uparrow \frac12.\;}
\]
特别地，\(I\) 在右端点具有垂直切线，而该端点不是解析边界点，而是对数型尖点。
\end{theorem}
\begin{proof}
记
\[
s(\alpha):=\inf_{\theta\in\RR}\bigl(P(\theta)-\alpha\theta\bigr).
\]
由 Legendre--Fenchel 对偶的归一化关系，
\[
I(\alpha)=P(0)-s(\alpha).
\]
另一方面，结论 \ref{con:xi-terminal-real-input-output-microcanonical-logcusp-second-order} 已给出：若记
\[
a:=\frac{d}{c}=2A,
\qquad
\delta:=\frac12-\alpha,
\]
则
\[
s\!\left(\frac12-\delta\right)
=
\log c-2\delta\log\delta+2\left(1+\log\frac{a}{2}\right)\delta+O(\delta^2).
\]
由于 \(a/2=A\)，并且推论 \ref{cor:real-input-40-rate-endpoints} 给出
\[
I\!\left(\frac12\right)=P(0)-\log c,
\]
故直接相减便得
\[
I(\alpha)
=
I\!\left(\frac12\right)
+2\varepsilon\log\varepsilon
-2(1+\log A)\varepsilon
+O(\varepsilon^2).
\]
第二个盒装式只是将 \(\log\varepsilon=-|\log\varepsilon|\) 代回后的重写。

再证导数公式。对每个 \(\alpha\in(0,\frac12)\)，由严格凸性可取唯一极值点 \(\theta(\alpha)\) 满足
\[
P'(\theta(\alpha))=\alpha,
\]
并有
\[
I'(\alpha)=\theta(\alpha).
\]
由定理 \ref{thm:xi-output-potential-zero-temp-square-root-critical-law}，
\[
\frac12-P'(\theta)=\frac{d}{2c}e^{-\theta/2}+O(e^{-\theta})
=
A e^{-\theta/2}\bigl(1+O(e^{-\theta/2})\bigr).
\]
于是当 \(\alpha\uparrow\frac12\) 时，
\[
\varepsilon
=
A e^{-\theta(\alpha)/2}\bigl(1+O(e^{-\theta(\alpha)/2})\bigr),
\]
从而
\[
e^{-\theta(\alpha)/2}=\frac{\varepsilon}{A}+O(\varepsilon^2).
\]
两边取对数得到
\[
\theta(\alpha)
=
-2\log\varepsilon+2\log A+O(\varepsilon),
\]
亦即
\[
I'(\alpha)
=
-2\log\!\Bigl(\frac12-\alpha\Bigr)+2\log A+O\!\Bigl(\frac12-\alpha\Bigr).
\]
因此 \(I'(\alpha)\to+\infty\)，右端点具有垂直切线；而 \(\varepsilon\log\varepsilon\) 项则表明该端点不可能解析。证毕。
\end{proof}

\begin{theorem}[输出位势右端点对偶参数的二阶反演与速率函数边界正规形]\label{thm:xi-output-potential-right-endpoint-second-order-inversion-normal-form}
\leanverified{paper\_xi\_output\_potential\_right\_endpoint\_second\_order\_inversion\_normal\_form}
沿用命题 \ref{prop:real-input-40-pressure-freezing}、推论 \ref{cor:real-input-40-rate-endpoints} 与结论 \ref{con:xi-terminal-real-input-output-microcanonical-logcusp-second-order} 的记号。令 \(\theta(\alpha)\) 为满足
\[
P'(\theta(\alpha))=\alpha
\]
的唯一对偶参数，并记
\[
\delta:=\frac12-\alpha\downarrow 0.
\]
则有二阶反演公式
\[
\boxed{\;
\theta\!\left(\frac12-\delta\right)
=
2\log\frac{d}{2c\,\delta}
+
\left(\frac{8ce}{d^2}-4\right)\delta
+O(\delta^2).\;}
\]
进一步，归一化速率函数满足
\[
\boxed{\;
I\!\left(\frac12-\delta\right)
=
\log\frac{\lambda(1)}{c}
-2\delta\log\frac{d}{2c\,\delta}
-2\delta
+\left(2-\frac{4ce}{d^2}\right)\delta^2
+O(\delta^3).\;}
\]
因此右端点的对数尖点不仅具有主系数 \(2\)，其二阶曲率修正也由 Puiseux 常数 \(c,d,e\) 完全显式决定。
\end{theorem}
\begin{proof}
记
\[
a:=\frac{d}{c},
\qquad
b:=\frac{e}{c}-\frac{d^2}{2c^2},
\qquad
x:=e^{-\theta/2}.
\]
由命题 \ref{prop:real-input-40-pressure-freezing}，
\[
P(\theta)
=
\frac{\theta}{2}+\log c+a x+b x^2+O(x^3),
\]
从而
\[
P'(\theta)
=
\frac12-\frac{a}{2}x-bx^2+O(x^3).
\]
于是极值条件 \(P'(\theta)=\frac12-\delta\) 等价于
\[
\delta=\frac{a}{2}x+bx^2+O(x^3).
\]
对该关系作幂级数反演可得
\[
x=\frac{2\delta}{a}-\frac{8b}{a^3}\delta^2+O(\delta^3).
\]
再取对数：
\[
\log x
=
\log\frac{2\delta}{a}-\frac{4b}{a^2}\delta+O(\delta^2),
\]
故
\[
\theta\!\left(\frac12-\delta\right)
=
-2\log x
=
2\log\frac{a}{2\delta}+\frac{8b}{a^2}\delta+O(\delta^2).
\]
将 \(a,b\) 代回，便得到
\[
\theta\!\left(\frac12-\delta\right)
=
2\log\frac{d}{2c\,\delta}
+
\left(\frac{8ce}{d^2}-4\right)\delta
+O(\delta^2).
\]

另一方面，结论 \ref{con:xi-terminal-real-input-output-microcanonical-logcusp-second-order} 已给出
\[
s\!\left(\frac12-\delta\right)
=
\log c
-2\delta\log\delta
+2\left(1+\log\frac{a}{2}\right)\delta
+\frac{4b}{a^2}\delta^2
+O(\delta^3).
\]
再由
\[
I(\alpha)=P(0)-s(\alpha),
\qquad
P(0)=\log\lambda(1),
\]
可得
\[
I\!\left(\frac12-\delta\right)
=
\log\frac{\lambda(1)}{c}
+2\delta\log\delta
-2\left(1+\log\frac{a}{2}\right)\delta
-\frac{4b}{a^2}\delta^2
+O(\delta^3).
\]
将其改写为
\[
I\!\left(\frac12-\delta\right)
=
\log\frac{\lambda(1)}{c}
-2\delta\log\frac{a}{2\delta}
-2\delta
-\frac{4b}{a^2}\delta^2
+O(\delta^3),
\]
再代回 \(a=d/c\) 与 \(b=e/c-d^2/(2c^2)\)，便得到
\[
I\!\left(\frac12-\delta\right)
=
\log\frac{\lambda(1)}{c}
-2\delta\log\frac{d}{2c\,\delta}
-2\delta
+\left(2-\frac{4ce}{d^2}\right)\delta^2
+O(\delta^3).
\]
证毕。
\end{proof}

\begin{theorem}[输出位势速率函数在左端点具有 Cram\'er--Lambert 型尖点]\label{thm:xi-output-potential-rate-left-endpoint-cramer-lambert}
沿用命题 \ref{prop:real-input-40-pressure-freezing} 与推论 \ref{cor:real-input-40-rate-endpoints} 的记号。则当 \(\alpha\downarrow 0\) 时，有
\[
\boxed{\;
I(\alpha)
=
I(0)+\alpha\log\alpha-(1+\log 4)\alpha+O(\alpha^2).\;}
\]
等价地，
\[
\boxed{\;
I(0)-I(\alpha)
=
\alpha|\log\alpha|+(1+\log 4)\alpha+O(\alpha^2).\;}
\]
并且
\[
\boxed{\;
I'(\alpha)=\log\alpha-\log4+O(\alpha)\to-\infty,\;}
\]
\[
\boxed{\;
I''(\alpha)=\frac1\alpha+O(1).\;}
\]
因此左端点同样不是解析端点，而是以主系数 \(1\) 控制的对数尖点。
\end{theorem}
\begin{proof}
由命题 \ref{prop:real-input-40-pressure-freezing} 在 \(u\to 0\) 时的展开，
\[
\lambda(u)=1+4u-5u^2-115u^3+2317u^4+O(u^5),
\]
故
\[
\log\lambda(u)=4u+O(u^2).
\]
另一方面，若记 \(\theta=\log u\)，则
\[
\alpha(u)=P'(\theta)=u\frac{\lambda'(u)}{\lambda(u)}=4u+O(u^2).
\]
因此可反解为
\[
u=\frac{\alpha}{4}+O(\alpha^2),
\qquad
\log u=\log\alpha-\log4+O(\alpha).
\]

再由 Legendre 对偶在内部点上的参数表示，
\[
I(\alpha(u))=\alpha(u)\log u-\log\lambda(u)+P(0).
\]
推论 \ref{cor:real-input-40-rate-endpoints} 给出
\[
I(0)=P(0)=\log\lambda(1)=\log(\varphi^2).
\]
并且由
\[
\log\lambda(u)=4u+O(u^2)=\alpha+O(\alpha^2)
\]
可得
\[
\begin{aligned}
I(\alpha)
&=
I(0)+\alpha\log u-\log\lambda(u)\\
&=
I(0)+\alpha(\log\alpha-\log4+O(\alpha))-\alpha+O(\alpha^2)\\
&=
I(0)+\alpha\log\alpha-(1+\log4)\alpha+O(\alpha^2).
\end{aligned}
\]
这便得到第一条与其等价重写。

对导数，仍记内部极值点为 \(\theta(\alpha)\)。则
\[
I'(\alpha)=\theta(\alpha)=\log u=\log\alpha-\log4+O(\alpha),
\]
从而 \(I'(\alpha)\to-\infty\)。再对上式求导，便得到
\[
I''(\alpha)=\frac1\alpha+O(1).
\]
因此左端点同样具有对数型非解析尖点。证毕。
\end{proof}

\begin{corollary}[输出位势速率函数两端尖点强度的刚性不对称]\label{cor:xi-output-potential-endpoint-log-cusp-rigidity-pair}
沿用定理 \ref{thm:xi-output-potential-rate-right-endpoint-log-cusp} 与定理 \ref{thm:xi-output-potential-rate-left-endpoint-cramer-lambert} 的记号。则有极限比值
\[
\boxed{\;
\lim_{\alpha\downarrow0}\frac{I(0)-I(\alpha)}{\alpha|\log\alpha|}=1,
\qquad
\lim_{\alpha\uparrow1/2}\frac{I(1/2)-I(\alpha)}{(1/2-\alpha)|\log(1/2-\alpha)|}=2.\;}
\]
因此，在本文固定的频率坐标 \(\alpha\) 中，两端对数尖点强度形成有序对
\[
\boxed{\;(1,2).\;}
\]
特别地，不存在端点交换局部对称
\[
I(0)-I(\alpha)=I\!\left(\frac12\right)-I\!\left(\frac12-\alpha\right)+o(\alpha|\log\alpha|)
\qquad
(\alpha\downarrow0),
\]
因而更不存在以
\[
\alpha\longmapsto \frac12-\alpha
\]
为主导阶模型的仿射端点反演对称。
\end{corollary}
\begin{proof}
第一条极限直接由定理 \ref{thm:xi-output-potential-rate-left-endpoint-cramer-lambert} 的主导项
\[
I(0)-I(\alpha)=\alpha|\log\alpha|+O(\alpha)
\]
得到。第二条极限则由定理 \ref{thm:xi-output-potential-rate-right-endpoint-log-cusp} 的主导项
\[
I\!\left(\frac12\right)-I(\alpha)
=
2\Bigl(\frac12-\alpha\Bigr)\Bigl|\log\!\Bigl(\frac12-\alpha\Bigr)\Bigr|
+O\!\left(\frac12-\alpha\right)
\]
立得。

若存在所述端点交换局部对称，则两侧在 \(\alpha|\log\alpha|\) 量级上的主导系数必须相同；但左端系数为 \(1\)，右端系数为 \(2\)，矛盾。故该局部对称不可能成立，而仿射反演 \(\alpha\mapsto\frac12-\alpha\) 只是其最直接的特例。证毕。
\end{proof}

\noindent 这几条结果进一步表明：输出位势的端点几何并非由单一对称冻结机制主宰。右端点同时受到零温四次 Perron 地基与 \(e^{-\theta/2}\) 尾部的控制，因此不仅产生主系数为 \(2\) 的对数尖点，而且其对偶反演与二阶曲率修正也被 \(c,d,e\) 完全显式锁定；左端点则由 \(u\to0\) 的 Cram\'er 型稀薄展开主导，仅留下主系数为 \(1\) 的对数尖点。于是，RH-sharp 五次阈值先于四次锚点出现的谱事实，与速率函数两端不可对称化的边界几何，便在同一条输出位势终端链上汇合起来。

\endinput


\begin{theorem}[最大密度层与平衡均值之间存在显式代数间隙而仍保留正熵]\label{thm:xi-real-input-ground-layer-mean-gap-positive-entropy}
\leanverified{paper\_xi\_real\_input\_ground\_layer\_mean\_gap\_positive\_entropy}
沿用推论 \ref{cor:real-input-40-alpha-max}、推论 \ref{cor:real-input-40-ground-entropy} 与表 \ref{tab:real_input_40_output_potential_cumulants_closed} 的记号，记
\[
\alpha_{\max}:=\frac12,
\qquad
\alpha_0:=P'(0)=\frac{23}{20}-\frac{7\sqrt5}{20},
\]
\[
\Delta_\mu:=\alpha_{\max}-\alpha_0
=
\frac{7\sqrt5-13}{20}
\approx 0.132623792124926.
\]
则对每个 \(n\) 及每条 \(\gamma\in\mathcal G_n\)，都有
\[
\frac{k_{\max}(n)}{n}
=
\frac12+O\!\left(\frac{1}{n}\right)
=
\alpha_0+\Delta_\mu+O\!\left(\frac{1}{n}\right).
\]
并且极端饱和层的总数满足
\[
|\mathcal G_n|
=
\exp\bigl((\log c)n+o(n)\bigr).
\]
因此在距离平衡均值 \(\Delta_\mu n+O(1)\) 的线性偏差层上，仍承载指数多的可行路径。
\end{theorem}
\begin{proof}
由推论 \ref{cor:real-input-40-ground-entropy}，
\[
k_{\max}(n)=\lfloor n/2\rfloor,
\qquad
|\mathcal G_n|=A_{n,k_{\max}(n)}
=
\exp\bigl((\log c)n+o(n)\bigr).
\]
因此
\[
\frac{k_{\max}(n)}{n}=\frac12+O\!\left(\frac{1}{n}\right).
\]
另一方面，表 \ref{tab:real_input_40_output_potential_cumulants_closed} 给出
\[
P'(0)=\frac{23}{20}-\frac{7\sqrt5}{20},
\]
故
\[
\frac12-P'(0)=\frac{7\sqrt5-13}{20}=\Delta_\mu>0.
\]
将两式合并，即得
\[
\frac{k_{\max}(n)}{n}
=
P'(0)+\Delta_\mu+O\!\left(\frac{1}{n}\right),
\]
并且同一偏差层的路径总数仍具有指数增长率 \(\log c\)。证毕。
\end{proof}

\paragraph{bin-fold 饱和窗口、零温微正则尖点与近地基壳层上界}
在 bin-fold 最大纤维的精确饱和判据、输出位势右端点的对数尖点以及地基层的正熵增长率都已固定之后，还可进一步把离散最大纤维、零温 Puiseux 展开与近端壳层计数压缩为同一条显式链。其共同特征在于：饱和窗口只由单个 Fibonacci 差值决定，而同一差值又经 Legendre--Fenchel 变换传递为右端点的对数尖点与近地基壳层的指数上界。

\leanverified{paper\_xi\_foldbin\_maxfiber\_saturation\_gap\_window\_closed}
\begin{theorem}[最大纤维饱和的差窗闭式与最大尾和公式]\label{thm:xi-foldbin-maxfiber-saturation-gap-window-closed}
沿用定义 \ref{def:K-of-m}、命题 \ref{prop:fold-bin-digit-dp} 与推论 \ref{cor:fold-bin-saturation} 的记号。对任意 \(m\ge 1\)，令
\[
K:=K(m),
\qquad
L:=K-m,
\qquad
d_{\max}^{\mathrm{bin}}(m):=d_m^{\mathrm{bin}}(0^m).
\]
再定义
\[
S_{\max}(m):=
\max\Bigl\{
\sum_{k=m+1}^{K}c_kF_{k+1}:\ c_k\in\{0,1\},\ c_kc_{k+1}=0
\Bigr\},
\]
并记
\[
\varepsilon_m:=
\begin{cases}
1,& L\equiv 0\pmod 2,\\
0,& L\equiv 1\pmod 2.
\end{cases}
\]
则 \(S_{\max}(m)\) 具有统一闭式
\[
\boxed{\;
S_{\max}(m)
=
\sum_{j=0}^{\lfloor (L-1)/2\rfloor}F_{K+1-2j}
=
F_{K+2}-F_{m+1+\varepsilon_m},
\;}
\]
其中当 \(L=0\) 时上式中的有限和按空和解释为 \(0\)。特别地，
\[
\boxed{\;
d_{\max}^{\mathrm{bin}}(m)=F_{L+2}
\iff
S_{\max}(m)\le 2^m-1
\iff
0\le F_{K+2}-2^m\le F_{m+1+\varepsilon_m}-1.
\;}
\]
\end{theorem}
\begin{proof}
由推论 \ref{cor:fold-bin-saturation}，最大纤维由 \(0^m\) 取得，并且
\[
d_{\max}^{\mathrm{bin}}(m)=F_{L+2}
\iff
S_{\max}(m)\le 2^m-1.
\]
因此只需把 \(S_{\max}(m)\) 写成统一闭式。

当 \(L=0\) 时尾部为空，故 \(S_{\max}(m)=0\)，而此时 \(K=m\) 且
\[
F_{K+2}-F_{m+1+\varepsilon_m}=F_{m+2}-F_{m+2}=0.
\]
以下设 \(L\ge 1\)。

由于权重 \(F_{k+1}\) 随 \(k\) 严格递增，若某个 admissible 尾部在位置 \(K\) 未取 \(1\)，则分两种情形：若 \(c_{K-1}=0\)，直接把 \(c_K\) 改为 \(1\) 即使尾和严格增大；若 \(c_{K-1}=1\)，则把 \(c_{K-1}\) 改为 \(0\) 并把 \(c_K\) 改为 \(1\)，尾和仍然严格增大。故任一极大尾部都必满足 \(c_K=1\)。去掉第 \(K\) 位后，同一论证可递归施加到区间 \(\{m+1,\dots,K-2\}\)，故极大模式只能是
\[
c_K=c_{K-2}=c_{K-4}=\cdots=1,
\]
其余各位皆为 \(0\)。因此
\[
S_{\max}(m)=\sum_{j=0}^{\lfloor (L-1)/2\rfloor}F_{K+1-2j}.
\]

再用恒等式
\[
\sum_{j=0}^{r}F_{n-2j}=F_{n+1}-F_{n-2r-1},
\]
取 \(n=K+1\)、\(r=\lfloor (L-1)/2\rfloor\)，得到
\[
S_{\max}(m)
=
F_{K+2}-F_{K-2\lfloor (L-1)/2\rfloor}.
\]
若 \(L=2s\) 为偶数，则 \(\lfloor (L-1)/2\rfloor=s-1\)，减项为 \(F_{m+2}\)；若 \(L=2s+1\) 为奇数，则减项为 \(F_{m+1}\)。这正好统一为
\[
S_{\max}(m)=F_{K+2}-F_{m+1+\varepsilon_m}.
\]

将该闭式代回饱和判据，便得到
\[
F_{K+2}-F_{m+1+\varepsilon_m}\le 2^m-1
\iff
F_{K+2}-2^m\le F_{m+1+\varepsilon_m}-1.
\]
而由 \(F_{K+1}\le 2^m-1<F_{K+2}\) 可知 \(2^m\le F_{K+2}\)，故左侧差值自动非负。证毕。
\end{proof}

\begin{theorem}[饱和点的指数级 Diophantine 逼近与有限性]\label{thm:xi-foldbin-saturation-diophantine-finiteness}
定义饱和点集合
\[
\mathcal S
:=
\Bigl\{
m\ge 1:\ d_{\max}^{\mathrm{bin}}(m)=F_{K(m)-m+2}
\Bigr\}.
\]
则存在常数 \(C>0\)，使得对每个 \(m\in\mathcal S\) 都有
\[
\boxed{\;
\Bigl|(K(m)+2)\log\varphi-m\log 2-\log\sqrt5\Bigr|
\le
C\Bigl(\frac{\varphi}{2}\Bigr)^m.
\;}
\]
因此 \(\mathcal S\) 只能是有限集；特别地，它具有零自然密度。
\end{theorem}
\begin{proof}
若 \(m\in\mathcal S\)，则由定理 \ref{thm:xi-foldbin-maxfiber-saturation-gap-window-closed}，
\[
0\le F_{K+2}-2^m\le F_{m+1+\varepsilon_m}-1\le F_{m+2}.
\]
令
\[
\widehat\varphi:=-\varphi^{-1},
\qquad
\Delta_m:=(K+2)\log\varphi-m\log 2-\log\sqrt5.
\]
把 Binet 公式
\[
F_n=\frac{\varphi^n-\widehat\varphi^n}{\sqrt5}
\]
代入，得到
\[
\left|
\frac{\varphi^{K+2}}{\sqrt5}-2^m
\right|
\le
F_{m+2}+\frac{|\widehat\varphi|^{K+2}}{\sqrt5}.
\]
由 \(F_{K+1}\le 2^m-1<F_{K+2}\) 可知 \(K=(\log 2/\log\varphi)m+O(1)\)，故
\[
F_{m+2}+\frac{|\widehat\varphi|^{K+2}}{\sqrt5}
=
O(\varphi^m).
\]
两边除以 \(2^m\) 即得
\[
|e^{\Delta_m}-1|
\le
C_1\Bigl(\frac{\varphi}{2}\Bigr)^m.
\]
由于 \(\varphi/2<1\)，当 \(m\) 足够大时右侧小于 \(1/2\)，于是
\[
|\Delta_m|
\le
2|e^{\Delta_m}-1|
\le
2C_1\Bigl(\frac{\varphi}{2}\Bigr)^m.
\]
把有限多个小 \(m\) 吸收到常数中，即得所示不等式。

另一方面，\(\Delta_m\) 是关于代数数 \(2\) 与 \(\varphi\) 的非零对数线性形式；由 Baker--W\"ustholz/Matveev 型下界，它在指数 \(m\) 上只能达到至多多项式级的小量，不可能对无穷多个 \(m\) 满足上述指数级估计。故 \(\mathcal S\) 必为有限集。这正是注 \ref{rem:fold-bin-finite} 所记录的机制。有限集当然具有零自然密度。证毕。
\end{proof}

\begin{theorem}[零温端点微正则包络的极小点与显式尖点闭式]\label{thm:xi-output-microcanonical-logcusp-legendre-closed}
沿用命题 \ref{prop:real-input-40-pressure-freezing} 的记号，定义
\[
h(\alpha):=\inf_{\theta\ge 0}\bigl(P(\theta)-\alpha\theta\bigr),
\qquad
0\le \alpha\le \frac12.
\]
再记
\[
a:=\frac{d}{c}>0,
\qquad
b:=\frac{e}{c}-\frac{d^2}{2c^2}.
\]
则存在 \(\delta_0>0\)，使得当
\[
\delta:=\frac12-\alpha\in(0,\delta_0)
\]
时，\(h(\alpha)\) 在 \([0,\infty)\) 上的极小点唯一，记为 \(\theta(\alpha)\)。并且
\[
\boxed{\;
\theta(\alpha)
=
-2\log\frac{2\delta}{a}
+\frac{8b}{a^2}\delta
+O(\delta^2),
\qquad
\delta\downarrow 0,
\;}
\]
以及
\[
\boxed{\;
h\!\left(\frac12-\delta\right)
=
\log c
-2\delta\log\delta
+2\left(1+\log\frac{a}{2}\right)\delta
+\frac{4b}{a^2}\delta^2
+O(\delta^3).
\;}
\]
\end{theorem}
\begin{proof}
由推论 \ref{cor:real-input-40-alpha-max} 与命题 \ref{prop:real-input-40-output-cumulants-closed}，
\[
P'(0)<\frac12.
\]
因此当 \(\delta>0\) 充分小时，极小点不可能停在边界 \(\theta=0\)，而必须满足驻点方程
\[
P'(\theta)=\frac12-\delta.
\]
再由命题 \ref{prop:real-input-40-pressure-freezing}，
\[
P(\theta)
=
\frac{\theta}{2}+\log c+a e^{-\theta/2}+b e^{-\theta}+O(e^{-3\theta/2}),
\]
从而
\[
P'(\theta)
=
\frac12-\frac{a}{2}e^{-\theta/2}-b e^{-\theta}+O(e^{-3\theta/2}).
\]
令
\[
u:=e^{-\theta/2},
\]
则驻点方程改写为
\[
\delta=\frac{a}{2}u+b u^2+O(u^3).
\]
对该幂级数作反演，得到
\[
u=\frac{2\delta}{a}-\frac{8b}{a^3}\delta^2+O(\delta^3).
\]
取对数并乘以 \(-2\)，便有
\[
\theta(\alpha)
=
-2\log u
=
-2\log\frac{2\delta}{a}
+\frac{8b}{a^2}\delta
+O(\delta^2).
\]
由于同一邻域内
\[
P''(\theta)=\frac{a}{4}e^{-\theta/2}+O(e^{-\theta})>0,
\]
故该驻点唯一，并且就是 \(h(\alpha)\) 的唯一极小点。

第二个展开正是结论 \ref{con:xi-terminal-real-input-output-microcanonical-logcusp-second-order} 的内容；将上式得到的 \(u\)-反演代回
\[
h(\alpha)=P(\theta(\alpha))-\alpha\theta(\alpha)
\]
亦可直接复得
\[
h\!\left(\frac12-\delta\right)
=
\log c
-2\delta\log\delta
+2\left(1+\log\frac{a}{2}\right)\delta
+\frac{4b}{a^2}\delta^2
+O(\delta^3).
\]
证毕。
\end{proof}

\begin{corollary}[零温端点微正则包络的斜率与曲率同时发散]\label{cor:xi-output-microcanonical-endpoint-slope-curvature-diverge}
沿用定理 \ref{thm:xi-output-microcanonical-logcusp-legendre-closed} 的记号。当 \(\alpha\uparrow 1/2\) 时，有
\[
\boxed{\;
h'(\alpha)
=
2\log\!\Bigl(\frac{2(1/2-\alpha)}{a}\Bigr)
+O\!\bigl(1/2-\alpha\bigr)
\longrightarrow -\infty,
\;}
\]
以及
\[
\boxed{\;
h''(\alpha)
=
-\frac{2}{1/2-\alpha}+O(1)
\longrightarrow -\infty.
\;}
\]
因此 \(h\) 在右端点不可能具有有限斜率的 \(C^1\) 延拓，更不存在有限曲率的二阶延拓。
\end{corollary}
\begin{proof}
令 \(\delta:=1/2-\alpha\)。由定理 \ref{thm:xi-output-microcanonical-logcusp-legendre-closed}，
\[
h\!\left(\frac12-\delta\right)
=
\log c
-2\delta\log\delta
+2\left(1+\log\frac{a}{2}\right)\delta
+\frac{4b}{a^2}\delta^2
+O(\delta^3).
\]
对 \(\delta\) 求导得
\[
\frac{d}{d\delta}
h\!\left(\frac12-\delta\right)
=
-2\log\delta
+2\log\frac{a}{2}
+\frac{8b}{a^2}\delta
+O(\delta^2).
\]
而 \(h'(\alpha)=-(d/d\delta)h(1/2-\delta)\)，故
\[
h'(\alpha)
=
2\log\frac{2\delta}{a}+O(\delta),
\]
这就给出第一式。再对上式求导，得到
\[
h''(\alpha)
=
-\frac{2}{\delta}+O(1)
=
-\frac{2}{1/2-\alpha}+O(1),
\]
因而二者都趋于 \(-\infty\)。证毕。
\end{proof}

\begin{theorem}[近地基壳层计数的统一 Legendre 上界]\label{thm:xi-near-ground-shell-counting-legendre-upper-bound}
沿用命题 \ref{prop:real-input-40-pressure-freezing} 与定理 \ref{thm:xi-real-input-40-output-edgeworth-quasipowers} 的记号。记
\[
Z_n(u)=\sum_{k=0}^{n}A_{n,k}u^k,
\qquad
k_n(\alpha):=\lfloor \alpha n\rfloor
\qquad
\bigl(0\le \alpha\le 1/2\bigr).
\]
则对任意固定的 \(\alpha\in[0,1/2]\)，都有
\[
\boxed{\;
\limsup_{n\to\infty}\frac{1}{n}\log A_{n,k_n(\alpha)}
\le
h(\alpha),
\;}
\]
其中 \(h\) 为定理 \ref{thm:xi-output-microcanonical-logcusp-legendre-closed} 中的微正则包络。特别地，当 \(\delta\downarrow 0\) 且 \(\alpha=1/2-\delta\) 时，
\[
\boxed{\;
\limsup_{n\to\infty}\frac{1}{n}\log A_{n,\lfloor(1/2-\delta)n\rfloor}
\le
\log c
-2\delta\log\delta
+2\left(1+\log\frac{a}{2}\right)\delta
+\frac{4b}{a^2}\delta^2
+O(\delta^3).
\;}
\]
\end{theorem}
\begin{proof}
任取 \(\theta\ge 0\)。由
\[
Z_n(e^\theta)=\sum_{k=0}^{n}A_{n,k}e^{\theta k},
\]
可得
\[
A_{n,k_n(\alpha)}e^{\theta k_n(\alpha)}
\le
Z_n(e^\theta).
\]
两边取对数并除以 \(n\)，注意到 \(k_n(\alpha)=\alpha n+O(1)\)，得到
\[
\frac{1}{n}\log A_{n,k_n(\alpha)}
\le
\frac{1}{n}\log Z_n(e^\theta)-\alpha\theta+O\!\left(\frac{1}{n}\right).
\]
对 \(n\to\infty\) 取上极限，并利用 Perron--Frobenius 压力极限
\[
\lim_{n\to\infty}\frac{1}{n}\log Z_n(e^\theta)=P(\theta),
\]
便有
\[
\limsup_{n\to\infty}\frac{1}{n}\log A_{n,k_n(\alpha)}
\le
P(\theta)-\alpha\theta.
\]
再对全部 \(\theta\ge 0\) 取下确界，便得到
\[
\limsup_{n\to\infty}\frac{1}{n}\log A_{n,k_n(\alpha)}
\le
\inf_{\theta\ge 0}\bigl(P(\theta)-\alpha\theta\bigr)
=
h(\alpha).
\]
当 \(\alpha=1/2-\delta\) 且 \(\delta\downarrow 0\) 时，只需把定理
\ref{thm:xi-output-microcanonical-logcusp-legendre-closed} 的端点展开代入即可。证毕。
\end{proof}

\noindent
综上，\(\log c\) 同时控制极端饱和层的最小完备外置率、无偏采样等待代价与微正则端点值；\(\Delta_\mu\) 继续刻画最大密度层与平衡均值之间的线性代数距离，而不消灭其正的指数熵；\(u_R\) 与 \(c^2=\rho(B_\infty)\) 之间的谱缝则说明 RH/Ramanujan 失效严格先于零温四次 Perron 尺度。进一步，bin-fold 最大纤维的离散异常饱和被压缩为单个 Fibonacci 差窗与指数级 Diophantine 逼近，而输出位势的零温 Puiseux 展开又把这一离散窗口传递为微正则端点的对数尖点与近地基壳层的统一上界。因而本段的零温结构可概括为
\[
\text{最大纤维饱和判据}
\Longrightarrow
\text{指数级 Diophantine 逼近}
\Longrightarrow
\text{零温端点 Puiseux 几何}
\Longrightarrow
\text{微正则对数尖点}
\Longrightarrow
\text{近地基壳层指数上界}.
\]
其中，从极端饱和层到其线性邻壳，主导几何始终由同一 Puiseux--Legendre 机制锁定。

\endinput

